\documentclass[10pt,a4paper]{article}

\usepackage[top=18mm,bottom=18mm,left=19mm,right=19mm]{geometry}
\usepackage{fontspec}
\usepackage{xeCJK}
\usepackage{titlesec}
\usepackage{enumitem}

\setCJKmainfont{Noto Serif CJK JP}
\setCJKsansfont{Noto Sans CJK JP}
\setmainfont{TeX Gyre Termes}
\setsansfont{TeX Gyre Heros}
\setlength{\parindent}{1em}
\setlength{\parskip}{0.25em}
\linespread{1.04}
\pagestyle{empty}
\titleformat{\section}{\sffamily\bfseries\normalsize}{}{0pt}{}
\titlespacing*{\section}{0pt}{0.55em}{0.15em}

\begin{document}

\begin{center}
  {\Large\sffamily\bfseries 研究計画書}\par
  \vspace{0.45em}
  {\large\sffamily\bfseries
  Auctionと討論型マルチLLMを組み合わせたGPU資源配分}
\end{center}

\section{研究背景}

大規模言語モデルの学習や科学技術計算では、多数の計算処理が、複数のGPUを共同で利用する。しかし、利用できるGPUの台数には限りがあり、すべての計算処理に希望どおりの台数を同時に与えることはできない。特に、計算の依頼が集中する時間帯には、ある処理へGPUを多く割り当てるほど、ほかの処理は開始できず、待ち時間や期限超過が増える。一方で、後から来る依頼に備えてGPUを空けすぎると、高価なGPUが使われないまま残る。このため、限られたGPUを、どの計算処理へ、何台、いつ割り当てるかを適切に決める必要がある。

この割り当てが難しい理由は、計算処理ごとにGPUを必要とする度合いが異なり、その状況も時間とともに変化するためである。期限が迫った処理や、その完了をほかの処理が待っている場合には、優先して実行する必要がある。その一方で、GPUを追加しても処理速度があまり向上しない計算もある。また、実行中の処理からGPUを頻繁に回収して別の処理へ割り当てると、その切替えにも時間と費用がかかる。したがって、希望台数や受付順だけで一律に決めるのではなく、GPUを追加する効果、処理の緊急性、全体の混み具合、切替えの負担をあわせて考える必要がある。

さらに、実際の運用では、数値だけでは表せない事情も割り当ての判断に影響する。たとえば、機器の故障や点検による制限、計算処理どうしの関係、依頼の取消し、管理者からの特別な優先指示などは、備考や通知の文章で伝えられる。このような情報を無視すれば重要な例外に対応できないが、すべての依頼について文章を詳しく分析すると、判断に時間と費用がかかる。そこで本研究では、通常の数値情報だけで判断できる計算処理は素早く扱い、文章による特別な事情がある計算処理だけを個別に検討することで、効率のよいGPU利用と、例外への柔軟な対応の両立を目指す。

\section{LLMを用いた先行研究}

近年、文章の内容を理解し、複数の条件を踏まえて案を作成できる大規模言語モデル（LLM）が注目されている。LLMは、決められた数値だけでなく、備考や通知に書かれた事情も読み取れるため、GPUの割り当てにも利用され始めている。

先行研究「LLM-Driven Adaptive Cloud Resource Scheduling」では、一つのLLMが、利用可能なGPUの台数、各計算処理の希望台数、優先度、期限、処理どうしの関係などをまとめて確認し、割り当て案を作成する。その後、数式に基づく最適化手法を用いて、GPUの総数を超えていないか、処理の順序や配置の条件に違反していないかを確認し、必要に応じて案を修正する。つまり、LLMが状況に応じた案を考え、最後に決められた規則で実行可能な形へ直す構成である。これにより、あらかじめすべての状況を規則として書き出すことが難しい場合でも、柔軟に割り当て案を作れることが示されている。

\section{先行研究に残された課題}

先行研究に残された中心的な課題は、LLMが生成した割り当て案の信頼性が低いことである。LLMは複数の条件を考慮して柔軟な案を作れる一方で、利用可能なGPUの台数を超えて割り当てたり、処理の順序や配置に関する条件に反した案を生成したりすることがある。実験では、守るべき条件をLLMへの指示に明記することで、条件に反する案の割合を34\%から8\%まで減らしている。しかし、指示を工夫しても条件違反を完全には防げておらず、LLMが作った案をそのまま実行することはできない。

先行研究では、この問題に対応するため、LLMが作った案を数式に基づく手法で確認し、条件に合うように修正している。この修正を行わず、LLMの案だけを用いた実験では、期限までに処理を終えられなかった割合が78.3\%に達した。この結果は、最終的な割り当ての信頼性をLLMだけでは確保できず、別の仕組みによる確認と修正が必要であることを示している。したがって、LLMの柔軟な判断を活用しながら、生成された案の信頼性をどのように高めるかが課題となる。

\section{提案手法}

本研究では、投入されたジョブを、運用上のテキスト情報の有無によって二つの経路に分ける。備考、障害・保守通知、依存関係、取消し、特別な優先指示などの\textbf{テキストがないジョブはauctionで処理する}。一方、このような\textbf{テキストがあるジョブだけを討論型マルチLLMで処理する}。これにより、通常の数値情報だけで判断できる多数のジョブではLLMを起動せず、推論時間と計算費用を削減する。

\textbf{1．テキストのないジョブを処理するauction}

Auctionは、各ジョブが追加のGPUをどの程度必要としているかを入札値で表し、クラスタ側がGPUを提供するための価格を提示して、両者が一致する範囲まで資源を配分する決定論的な仕組みである。自然言語の生成を行わず、数値計算だけで結果を決めるため、高速であり、同じ入力には常に同じ結果を返す。

まず、各ジョブは1台目、2台目、3台目というように、追加GPUごとの限界価値を入札する。入札値には、そのGPUによって増える有効な処理量、期限に遅れる危険性、待機時間、および実行中のジョブを拡大・縮小する費用を反映する。一般にGPUを追加するほど得られる効果は小さくなるため、各ジョブの入札値はGPU台数の増加に伴って低下する曲線とする。これにより、単に要求台数が多いジョブではなく、追加GPUから大きな効果を得られるジョブが優先される。

次に、供給側は、クラスタの空きGPU数、現在の混雑度、資源の断片化、待機中ジョブ数、および今後到着する重要ジョブの可能性から、1台ごとの売値を計算する。多くのGPUを配分して残りが少なくなるほど売値を高くすることで、資源を一つのジョブへ過剰に配分せず、将来の需要に備える。

市場を清算する際には、全ジョブの入札値を高い順に並べ、供給側の売値と比較する。$q$台目について入札値が売値以上であれば、そのGPUを該当ジョブへ配分する。入札値が売値を下回った時点で配分を終了する。すなわち、
\[
Q=\max\{q\mid b_{(q)}\geq a_q\}
\]
として配分台数$Q$を決定する。ここで$b_{(q)}$は高い順に並べた$q$番目の入札値、$a_q$は$q$台目の売値である。同じ値の場合の順序もあらかじめ固定し、結果を再現可能にする。さらに、現在の配分からGPUを移動する利益がジョブの縮小・再配置費用を上回る場合だけ変更を適用し、頻繁な配分変更を防止する。最後に、総配分数が容量を超えていないか、負の配分や未知のジョブがないかを決定論的に検証する。

\textbf{2．テキストのあるジョブを処理する討論型マルチLLM}

テキストが付与されたジョブでは、数値だけでは例外の意味を判断できないため、需要側LLM、供給側LLM、審判LLMを起動する。需要側LLMは、期限、依存関係、取消し、障害の影響などからジョブがGPUを必要とする理由を述べる。供給側LLMは、クラスタ全体の容量、公平性、他ジョブへの影響、保守上の制約から反対意見または修正案を示す。両者は互いの主張を確認して再判断し、審判LLMがauctionによる通常ジョブの配分と残り容量も考慮して最終案を決定する。

最後に、auctionの結果と討論の結果を統合し、全ジョブの総配分がクラスタ容量を超えないことを検証する。討論結果が制約に違反した場合、または制限時間内に得られなかった場合は、そのジョブの数値情報を用いたauctionの結果へ戻す。この構成により、通常ジョブはauctionで高速に処理し、自然言語による例外判断が必要なジョブだけにLLMの推論能力を用いる。

\end{document}
